% N2 --- Thevenin e Norton: 11 questoes com topologias distintas
% Revisao visual: portas roteadas fora dos ramos e maior separacao de rotulos.
\blocoaplicacao

\begin{questao}[2]{}
Determine o equivalente de Thévenin visto de A--B.
\circuitoq{\begin{circuitikz}[american,scale=.86,transform shape,line width=.8pt]
\coordinate (g) at (4.2,0); \coordinate (a) at (4.2,4.0);
\draw (0,0) to[\fonteV,l^={$\SI{24}{\volt}$}] (0,4.0) to[R,l={$\SI{6}{\kilo\ohm}$}] (a);
\draw (a) to[R,l_={$\SI{3}{\kilo\ohm}$}] (g)--(0,0);
\draw (a)--(6.6,4.0) node[ocirc]{} node[right]{A};
\draw (g)--(6.6,0) node[ocirc]{} node[right]{B};
\fill (a) circle (1.1pt); \node[ground] at(g){};
\end{circuitikz}}
\resposta{$V_{Th}=\SI{8}{\volt}$; $R_{Th}=\SI{2}{\kilo\ohm}$}
\resolucao{Divisor em vazio: $24\cdot3/(6+3)=8$ V. Com a fonte em curto, $6\parallel3=2$ k$\Omega$.}
\end{questao}

\begin{questao}[2]{}
Determine $V_{Th}$ e $R_{Th}$ vistos na porta à direita.
\circuitoq{\begin{circuitikz}[american,scale=.82,transform shape,line width=.8pt]
\coordinate (g) at (2.8,0); \coordinate (a) at (2.8,4.0); \coordinate (b) at (7.6,4.0);
\draw (0,0) to[isource,l^={$\SI{2}{\milli\ampere}$}] (0,4.0)--(a);
\draw (a) to[R,l_={$\SI{4}{\kilo\ohm}$}] (g)--(0,0);
\draw (a) to[R,l={$\SI{2}{\kilo\ohm}$}] (b) node[ocirc]{} node[right]{A};
\draw (g)--(7.6,0) node[ocirc]{} node[right]{B};
\fill (a) circle (1.1pt); \node[ground] at(g){};
\end{circuitikz}}
\resposta{$V_{Th}=\SI{8}{\volt}$; $R_{Th}=\SI{6}{\kilo\ohm}$}
\resolucao{Em aberto não circula corrente no resistor de 2 k$\Omega$, então $V_{Th}=2$ mA$\times4$ k$\Omega=8$ V. Com a fonte aberta, a porta vê $2+4=6$ k$\Omega$.}
\end{questao}

\begin{questao}[2]{}
Duas fontes reais aterradas alimentam a mesma porta. Determine o equivalente.
\circuitoq{\begin{circuitikz}[american,scale=.68,transform shape,line width=.8pt]
\coordinate (g) at (5.4,0); \coordinate (a) at (5.4,4.4);
\coordinate (pa) at (7.2,5.8); \coordinate (pb) at (7.2,-1.3);
\draw (0,0) to[\fonteV,l^={$\SI{20}{\volt}$}] (0,4.4) to[R,l={$\SI{5}{\kilo\ohm}$}] (a);
\draw (11.0,0) to[\fonteV,l_={$\SI{10}{\volt}$}] (11.0,4.4) to[R,l_={$\SI{10}{\kilo\ohm}$}] (a);
\draw (0,0)--(g)--(11.0,0);
\draw (a)--(5.4,5.8)--(pa) node[ocirc]{} node[right]{A};
\draw (g)--(5.4,-1.3)--(pb) node[ocirc]{} node[right]{B};
\fill (a) circle (1.1pt); \node[ground] at(g){};
\end{circuitikz}}
\resposta{$V_{Th}\approx\SI{16.67}{\volt}$; $R_{Th}\approx\SI{3.33}{\kilo\ohm}$}
\resolucao{$V_{Th}=(20/5+10/10)/(1/5+1/10)=16{,}67$ V. Zerando as fontes, $R_{Th}=5\parallel10=3{,}33$ k$\Omega$.}
\end{questao}

\begin{questao}[2]{}
A porta está no nó B. Determine seu equivalente de Thévenin.
\circuitoq{\begin{circuitikz}[american,scale=.74,transform shape,line width=.8pt]
\coordinate (g) at (3.8,0); \coordinate (a) at (3.8,4.2); \coordinate (b) at (9.2,4.2);
\draw (0,0) to[\fonteV,l^={$\SI{12}{\volt}$}] (0,4.2) to[R,l={$R_1=\SI{2}{\kilo\ohm}$}] (a);
\draw (a) to[R,l_={$R_2=\SI{4}{\kilo\ohm}$}] (g)--(0,0);
\draw (a) to[R,l={$R_3=\SI{3}{\kilo\ohm}$}] (b) node[ocirc]{} node[right]{B};
\draw (g)--(9.2,0) node[ocirc]{} node[right]{terra};
\fill (a) circle (1.1pt); \node[Vcol,fill=white,inner sep=1.4pt] at ($(a)+(0,.65)$) {$A$};
\node[ground] at(g){};
\end{circuitikz}}
\resposta{$V_{Th}=\SI{8}{\volt}$; $R_{Th}\approx\SI{4.33}{\kilo\ohm}$}
\resolucao{Em aberto, $R_3$ não conduz e $V_B=12\cdot4/(2+4)=8$ V. Com a fonte zerada, $R_{Th}=R_3+(R_1\parallel R_2)=3+4/3=4{,}33$ k$\Omega$.}
\end{questao}

\begin{questao}[2]{}
Determine o equivalente de Thévenin da fonte de Norton com resistor série até a porta.
\circuitoq{\begin{circuitikz}[american,scale=.82,transform shape,line width=.8pt]
\coordinate (g) at (3.0,0); \coordinate (a) at (3.0,4.0); \coordinate (b) at (7.8,4.0);
\draw (0,0) to[isource,l^={$\SI{5}{\milli\ampere}$}] (0,4.0)--(a);
\draw (a) to[R,l_={$\SI{3}{\kilo\ohm}$}] (g)--(0,0);
\draw (a) to[R,l={$\SI{2}{\kilo\ohm}$}] (b) node[ocirc]{} node[right]{A};
\draw (g)--(7.8,0) node[ocirc]{} node[right]{B};
\fill (a) circle (1.1pt); \node[ground] at(g){};
\end{circuitikz}}
\resposta{$V_{Th}=\SI{15}{\volt}$; $R_{Th}=\SI{5}{\kilo\ohm}$}
\resolucao{Em vazio o resistor série de 2 k$\Omega$ não conduz, então $V_{Th}=5$ mA$\times3$ k$\Omega=15$ V. Com a fonte aberta, $R_{Th}=3+2=5$ k$\Omega$.}
\end{questao}

\begin{questao}[2]{}
Determine $V_{Th}$ e $R_{Th}$ vistos em A--B.
\circuitoq{\begin{circuitikz}[american,scale=.70,transform shape,line width=.8pt]
\coordinate (g) at (4.8,0); \coordinate (a) at (4.8,4.4);
\coordinate (pa) at (9.8,5.7); \coordinate (pb) at (9.8,-1.3);
\draw (0,0) to[\fonteV,l^={$\SI{12}{\volt}$}] (0,4.4) to[R,l={$\SI{4}{\ohm}$}] (a);
\draw (a) to[R,l_={$\SI{6}{\ohm}$}] (g)--(0,0);
\draw (8.4,0) to[isource,l_={$\SI{2}{\ampere}$}] (8.4,4.4)--(a); \draw (8.4,0)--(g);
\draw (a)--(4.8,5.7)--(pa) node[ocirc]{} node[right]{A};
\draw (g)--(4.8,-1.3)--(pb) node[ocirc]{} node[right]{B};
\fill (a) circle (1.1pt); \node[ground] at(g){};
\end{circuitikz}}
\resposta{$V_{Th}=\SI{12}{\volt}$; $R_{Th}=\SI{2.4}{\ohm}$}
\resolucao{A condutância da porta é $1/4+1/6=5/12$ S. A fonte de 12 V equivale a 3 A e soma-se à fonte de 2 A: $V_{Th}=5/(5/12)=12$ V. Zerando as fontes, $R_{Th}=4\parallel6=2{,}4\,\Omega$.}
\end{questao}

\begin{questao}[2]{}
O equivalente abaixo alimentará três cargas distintas. Complete as correntes e identifique a condição de maior potência.
\circuitoq{\begin{circuitikz}[american,scale=.86,transform shape,line width=.8pt]
\draw (0,0) to[\fonteV,l^={$V_{Th}=\SI{9}{\volt}$}] (0,3.8)
      to[R,l={$R_{Th}=\SI{3}{\ohm}$}] (4.2,3.8)
      to[R,l={$R_L$}] (4.2,0)--(0,0);
\node[right,align=left,font=\footnotesize,draw=cinza,rounded corners,inner sep=3pt]
      at(6.3,1.9){$R_L=3\,\Omega$\\$R_L=6\,\Omega$\\$R_L=15\,\Omega$};
\end{circuitikz}}
\resposta{$I_L=\SI{1.5}{\ampere}$, $\SI{1}{\ampere}$ e $\SI{0.5}{\ampere}$; maior potência em $R_L=\SI{3}{\ohm}$}
\resolucao{$I_L=9/(3+R_L)$. As correntes são 1,5 A, 1 A e 0,5 A. A máxima transferência ocorre quando $R_L=R_{Th}=3\,\Omega$.}
\end{questao}

\begin{questao}[2]{}
Uma caixa-preta fornece $\SI{30}{\volt}$ em vazio e $\SI{24}{\volt}$ com $R_L=\SI{10}{\ohm}$. Determine o equivalente e $I_{sc}$.
\circuitoq{\begin{circuitikz}[american,scale=.84,transform shape,line width=.8pt]
\draw[dashed,cinza,line width=.8pt,rounded corners=6pt] (0,0) rectangle (3.4,3.9);
\node[cinza] at(1.7,2.0){caixa-preta};
\draw (3.4,3.1)--(5.0,3.1) coordinate(a); \draw (3.4,.8)--(5.0,.8) coordinate(b);
\draw (a) to[R,l={$R_L=\SI{10}{\ohm}$}] (b);
\draw[<->,Vcol,line width=.9pt] (6.0,1.0)--(6.0,2.9) node[midway,right] {$\SI{24}{\volt}$};
\node[above,font=\footnotesize,fill=white,inner sep=1.5pt] at(1.7,4.15){$V_{oc}=\SI{30}{\volt}$};
\end{circuitikz}}
\resposta{$V_{Th}=\SI{30}{\volt}$; $R_{Th}=\SI{2.5}{\ohm}$; $I_{sc}=\SI{12}{\ampere}$}
\resolucao{Com a carga, $I=24/10=2{,}4$ A e a queda interna é 6 V; logo $R_{Th}=6/2{,}4=2{,}5\,\Omega$. Então $I_{sc}=30/2{,}5=12$ A.}
\end{questao}

\begin{questao}[2]{}
A mesma rede é observada em duas portas. Determine os equivalentes em A e B.
\circuitoq{\begin{circuitikz}[american,scale=.68,transform shape,line width=.8pt]
\coordinate (g) at (4.0,0); \coordinate (a) at (4.0,4.5); \coordinate (b) at (10.0,4.5);
\coordinate (pa) at (6.0,6.0);
\draw (0,0) to[\fonteV,l^={$\SI{18}{\volt}$}] (0,4.5) to[R,l={$\SI{6}{\kilo\ohm}$}] (a);
\draw (a) to[R,l_={$\SI{3}{\kilo\ohm}$}] (g)--(0,0);
\draw (a) to[R,l={$\SI{4}{\kilo\ohm}$}] (b) node[ocirc]{} node[right]{B};
\draw (a)--(4.0,6.0)--(pa) node[ocirc]{} node[right]{A};
\draw (g)--(10.0,0) node[ocirc]{} node[right]{terra};
\fill (a) circle (1.1pt); \node[ground] at(g){};
\end{circuitikz}}
\resposta{Porta A: $V_{Th}=\SI{6}{\volt}$, $R_{Th}=\SI{2}{\kilo\ohm}$; porta B: $V_{Th}=\SI{6}{\volt}$, $R_{Th}=\SI{6}{\kilo\ohm}$}
\resolucao{Em vazio, nenhuma corrente passa no resistor até B, então ambas as portas têm 6 V. Em A, $R_{Th}=6\parallel3=2$ k$\Omega$; em B soma-se o resistor de 4 k$\Omega$, resultando 6 k$\Omega$.}
\end{questao}

\begin{questao}[2]{}
Determine o equivalente de Norton da rede.
\circuitoq{\begin{circuitikz}[american,scale=.84,transform shape,line width=.8pt]
\coordinate (g) at (4.2,0); \coordinate (a) at (4.2,4.0);
\draw (0,0) to[isource,l^={$\SI{4}{\ampere}$}] (0,4.0)--(a);
\draw (a) to[R,l_={$\SI{12}{\ohm}$}] (2.8,0)--(g);
\draw (a) to[R,l={$\SI{6}{\ohm}$}] (5.8,0)--(g); \draw (0,0)--(g);
\draw (a)--(7.8,4.0) node[ocirc]{} node[right]{A};
\draw (g)--(7.8,0) node[ocirc]{} node[right]{B};
\node[ground] at(g){};
\end{circuitikz}}
\resposta{$I_N=\SI{4}{\ampere}$; $R_N=\SI{4}{\ohm}$}
\resolucao{Em curto, toda a corrente da fonte vai para a porta, logo $I_N=4$ A. Com a fonte aberta, $R_N=12\parallel6=4\,\Omega$.}
\end{questao}

\begin{questao}[2]{}
Uma carga $R_L=\SI{5}{\ohm}$ é conectada ao equivalente de Norton mostrado. Determine $I_L$ e $U_L$.
\circuitoq{\begin{circuitikz}[american,scale=.86,transform shape,line width=.8pt]
\coordinate (g) at (4.2,0); \coordinate (a) at (4.2,4.0);
\draw (0,0) to[isource,l^={$I_N=\SI{3}{\ampere}$}] (0,4.0)--(a);
\draw (a) to[R,l_={$R_N=\SI{6}{\ohm}$}] (2.8,0)--(g);
\draw (a) to[R,l={$R_L=\SI{5}{\ohm}$},i>^={$I_L$}] (6.2,0)--(g);
\draw (0,0)--(g); \node[ground] at(g){};
\end{circuitikz}}
\resposta{$I_L\approx\SI{1.64}{\ampere}$; $U_L\approx\SI{8.18}{\volt}$}
\resolucao{Por divisão de corrente, $I_L=3\cdot6/(6+5)=18/11\approx1{,}64$ A. Então $U_L=5I_L\approx8{,}18$ V.}
\end{questao}